\section{Why this matters for AI agents}
\label{sec:agents}

The argument of this section is a single claim with four parts:

\begin{keyidea}
An AI agent's own knowledge is \textbf{not canonical, not addressable, not
verifiable, and not attributable}. Those four properties are exactly what
\odag{} provides, and they are provided \emph{because} of the refusals in
Section~\ref{sec:keep}, not in spite of them.

The corollary for design: a knowledge substrate for agents should sell
\emph{guarantees}, not \emph{expressiveness}. Reasoning is now abundant and
cheap --- every agent carries a flexible reasoner in its weights. Agreement is
scarce and expensive.
\end{keyidea}

\subsection{The four missing properties}

Take them one at a time, because each names a concrete failure mode.

\paragraph{Not canonical.} Ask a model the same question twice, phrased
differently, and you may get two answers that are both reasonable and
mutually inconsistent. There is no normal form for what a model knows, and
therefore no way to tell whether two models --- or the same model on Tuesday
--- know the same thing. \odag{}: equal knowledge has equal bytes (G1), so
``do we agree?'' is a string comparison.

\paragraph{Not addressable.} You cannot point at a fact in a model. There is
no reference you can put in a contract, cite in a document, or hand to a
counterparty that denotes ``the state of knowledge you were using''. \odag{}:
a root is a 32-byte reference to the entire knowledge state, immutable and
retrievable forever.

\paragraph{Not verifiable.} A model's output cannot be checked except by
another model or a human. This is the hallucination problem in its structural
form \citep{ji2023}: the failure is not that models are often wrong but that
being wrong is \emph{indistinguishable in the output} from being right.
\odag{}: an answer comes with a certificate a third party can check with
nothing but the root (\S\ref{sec:agents-verify}).

\paragraph{Not attributable.} When a model states something, there is no
answer to ``who says?'' beyond ``the model''. In a multi-agent system, where
one agent consumes another's output, this compounds silently. \odag{}: every
write is a signed record naming its author and the root they were looking at.

\subsection{This is not an argument against retrieval}

The obvious objection is that retrieval-augmented generation
\citep{lewis2020} already grounds models in external data. It does, and the
two approaches are complementary rather than competing. It is worth being
precise about the division of labour, because getting it wrong wastes both.

\begin{center}\small
\begin{tabularx}{\textwidth}{@{}p{3.5cm} X X@{}}
\toprule
& \textbf{Vector retrieval} & \textbf{\odag{} query}\\
\midrule
Input & vague, natural language & exact terms\\
Recall & good on paraphrase and synonymy & only what was filed\\
Precision & approximate top-$k$; no guarantee & exact: the complete
  intersection, or nothing\\
``Is this all of them?'' & unanswerable & yes, with a count\\
``Why this result?'' & similarity score & the path of asserted edges\\
Reproducible? & depends on index and model version & yes, given the root\\
Two systems agree? & incomparable spaces & compare 32 bytes\\
\bottomrule
\end{tabularx}
\end{center}

\begin{keyidea}
\textbf{The embedding proposes; the graph disposes.} Use vector search to go
from a vague human request to candidate terms; use \odag{} to answer exactly,
completely, and citably once the terms are fixed. The failure mode to avoid is
using approximate retrieval for a question whose answer must be complete ---
``list every refundable booking'' is not a similarity query, and a top-$k$
answer to it is wrong in a way nothing downstream can detect.
\end{keyidea}

\subsection{What the agent surface actually looks like}

\odag{} exposes itself to agents over the Model Context Protocol
\citep{mcp2024} as \cmd{odag-mcp}. Here is a real session against the running
example (transport framing removed, JSON reformatted):

\begin{lstlisting}
-> {"method":"tools/call","params":{"name":"about","arguments":{}}}

<- { "store": ".../store.od",
     "items": 37,
     "top": ["Booked","Japan","Refundable","Travel","dimension"],
     "dimensions": {"weight":"linear-dimension","time":"calendar-dimension",
                    "geo":"prefix-dimension", ...},
     "registry_version": "4.1",
     "surface_version": "0.1",
     "contract": "0.1",
     "capabilities": ["about","canon","describe","is_below",
                      "overlapping","query"],
     "root": "6f7342307f2fc4531c75f0df2465941f37822dd5f8cac5480e...",
     "annotations": {} }
\end{lstlisting}

Five design decisions are visible in that one response, and each is a
deliberate answer to one of the four missing properties.

\paragraph{1. Discoverability without downloading.} \cmd{about} lets an agent
learn what a store is \emph{about} --- its top-level categories, its declared
dimensions, its size --- before deciding whether to query it. Note that this
record is \emph{computed on demand and never stored}: a summary living inside
the store would make the root depend on a description of itself, since the
item count is part of the description and the description would be part of the
hash.

\paragraph{2. Every answer cites its root.} Not as decoration. It converts an
answer from a perishable claim into a permanent one: ``\texttt{get(Japan,
Refundable)} at root \texttt{6f734230\ldots} was these three documents'' is
true forever, checkable by anyone, and immune to the store changing
afterwards. An agent that caches answers with their roots has no
cache-invalidation problem, because it is not caching a guess about the
present; it is recording a fact about a past state.

\paragraph{3. The interpretation context is pinned.} \texttt{registry\_version}
and \texttt{contract} tell the agent which arithmetic and which guarantees
were in force. A verifier running different arithmetic must refuse rather than
misinterpret.

\paragraph{4. Answers are extensible objects, never bare lists.} Every
response carries a namespaced \texttt{annotations} map, empty today. New
information --- guarantee status, economic bonds, confidence from some other
system --- can arrive under its own namespace without changing any existing
answer shape, and readers ignore namespaces they do not know.

\paragraph{5. Canonical echo.} Every answer restates the terms in their
\emph{stored} form:

\begin{lstlisting}
-> {"name":"query","arguments":{"terms":["Japan","Refundable"]}}
<- { "terms": ["Japan","Refundable"],
     "items": ["kyoto-inn.pdf","outbound.pdf","rail-pass.pdf"],
     "count": 3, "truncated": false,
     "root": "6f7342307f2fc4531c75f0df2465941f37822dd5f8cac5480e...",
     "contract": "0.1", "annotations": {} }
\end{lstlisting}

With typed values the echo earns its keep: an agent that asks about
\verb|weight(400g)| is told the store calls that \verb|weight(2/5kg)|. It has
learned the canonical spelling, so its next assertion will collide correctly
with everyone else's rather than creating a near-duplicate. \emph{The canonical
echo is how an agent discovers the store's identity conventions without being
told them.}

Note also \texttt{count} beside \texttt{items}, and \texttt{truncated}. The
command-line interface truncates long output by default because a human at a
terminal can see that it did; the agent surface has \textbf{no default limit},
because an agent cannot see a terminal and so cannot see a truncation. When a
limit is set explicitly, \texttt{count} still reports the complete size. This
is a small thing that illustrates a general principle: \emph{an interface for
machines must make its own incompleteness machine-readable.}

\subsection{Closing the world by pinning, not by decree}
\label{sec:asof}

Agents constantly need non-monotone answers: ``is there anything under
\textsf{Japan} that is \emph{not} \textsf{Booked}?'', ``how many?'', ``does
this store contain X at all?''. Section~\ref{sec:mergeability} explained why
the living store cannot answer these: in a merging system the answer can
shrink when news arrives, so two replicas would disagree.

The escape is scoping rather than prohibition:

\begin{keyidea}
\textbf{Monotone questions may be asked of the living store. Non-monotone
questions must name a root.} Indexed by a root, negation, aggregation,
universal quantification and absence claims are pure, deterministic,
eternally recomputable facts. ``\texttt{get(A and not B)} at root
\texttt{21728cd9\ldots} $=$ \{\ldots\}'' is an immutable claim anyone can
verify by re-execution.
\end{keyidea}

The corollary for an agent is worth stating as a rule of practice: \emph{an
agent that needs a closed world does not ask the store to close. It closes the
world itself, by pinning a root.} Every tool on the surface takes an
\texttt{as\_of} parameter for exactly this.

\subsection{Fail-closed answers}

\cmd{is\_below} answers \emph{true} only when the graph plus exact arithmetic
witness it. \emph{False} means ``not derivable from what is here'', never
``unknown but plausible'' (G4).

This is the opposite of a language model's default behaviour, and it is what
makes the operation composable with a model. An agent can use the model for
the plausible and the store for the certain, and --- crucially --- can
\emph{tell which is which}. A system where the two are interchangeable is a
system where the guarantees of the exact half are worth nothing.

\subsection{Verification: from ``trust the store'' to ``trust nobody''}
\label{sec:agents-verify}

Because the persistence layer is a canonically encoded hash structure, ``this
record is in this version'' has a small proof: the exact stored bytes along
the one path where the record must live, checkable by recomputing hashes.
Unusually, \emph{absence} is provable too --- the encoding is canonical, so a
key has exactly one possible location, and exhibiting that location empty is a
proof it is not there.

\odag{} builds subsumption certificates on top. The design is
\emph{re-execution over authenticated fragments}: the prover records the real
\cmd{is\_below} walk, adds the order-invariant dependency closure (so a
verifier whose walk explores a different path is still covered), proves every
touched record, and the verifier re-runs the \emph{real} algorithm over a
store that serves only proof-verified records. Semantics stay single-sourced;
a coverage gap fails verification rather than validating a wrong answer.

\begin{measured}
Measured round trip (Experiment 9):
\begin{center}\ttfamily\footnotesize
\begin{tabular}{@{}ll@{}}
root & d53e10b9ee4a06edac177b23d40a3ddab9a8f9a890a4397ac16f50548acc39ea\\
certificate & ontodag-is-below-certificate v1, 4 records proved,
  registry 4.1\\
claim & outbound.pdf is below Travel $\to$ True\\
verify\_below(cert, root) & True \quad (no store, no network)\\
tampered (result flipped) & CertificateError\\
negative answer & outbound.pdf below Hotel $\to$ False, verified False\\
\end{tabular}
\end{center}
\end{measured}

The general pattern --- a data structure that ships proofs of its own answers
--- is what \citet{miller2014} call an authenticated data structure; what is
unusual here is that the \emph{semantic} operation, not merely record lookup,
is the thing being attested.

Both polarities work; a certificate may be incomplete (verification fails) but
never misleading (a wrong answer cannot verify). So an agent can hand another
agent a conclusion \emph{and the means to check it}, without handing over the
store. The answer travels; the data does not.

\begin{figure}[t]
\centering
\begin{tikzpicture}[font=\small,
  b/.style={draw, rounded corners=3pt, align=center, inner sep=5pt,
            minimum height=1.0cm, text width=3.0cm},
  ar/.style={-{Stealth[length=2.4mm]}, thick},
  lab/.style={font=\scriptsize\itshape, text=odgrey, align=center}]

\node[b, draw=odblue] (store) at (0,0)
  {\odag{} at root $R$\\[1pt]\scriptsize canonical, hashed};
\node[b, fill=odsand] (a1) at (4.6,0)
  {Agent A\\[1pt]\scriptsize holds the store};
\node[b, draw=odgreen] (cert) at (9.2,0)
  {answer $+$ certificate\\[1pt]\scriptsize
   \{format, root, proofs\}};
\node[b, fill=odsand] (a2) at (9.2,-2.4)
  {Agent B\\[1pt]\scriptsize holds only $R$};
\node[b, draw=odgreen, text width=4.3cm] (check) at (3.6,-2.4)
  {re-run the real \cmd{is\_below}\\[1pt]\scriptsize
   over proof-verified records};

\draw[ar] (store) -- node[lab, above]{query} (a1);
\draw[ar] (a1)    -- node[lab, above]{prove} (cert);
\draw[ar] (cert)  -- node[lab, right, xshift=1mm]{send} (a2);
\draw[ar] (a2)    -- node[lab, above]{verify} (check);
\draw[ar, odgreen, dashed] (check.west) to[out=180, in=180, looseness=1.1]
  node[lab, left, xshift=-1mm, text width=1.9cm]{true, or\\refuse} (store.west);

\node[lab, text width=12.5cm, anchor=north] at (5.0,-3.5)
  {Agent B needs no copy of the store, no network access, and no trust in
   Agent A --- only the 32 bytes of $R$ and the verifier code. The dashed
   return is a claim about $R$, not a request to the store.};
\end{tikzpicture}
\caption{Verification without disclosure. Because the root commits to the
whole state and the encoding is canonical, a subsumption answer can be
shipped with everything needed to check it.}
\label{fig:verify}
\end{figure}

\subsection{Writing: propose, echo, confirm}

Letting agents \emph{write} raises different problems. \odag{}'s write surface
answers three of them structurally.

\paragraph{Idempotence, by construction.} Re-filing the same fact is a graph
no-op: the same root comes out. An agent that retries after a timeout, or two
agents that independently reach the same conclusion, cannot corrupt anything.
This falls straight out of the canonical form --- it is not a feature that had
to be built.

\paragraph{Compare-and-confirm.} A write is two steps. \cmd{propose\_put}
returns the canonical echo (what the stored spellings will be, which claims
this implies, what is already true, which parents do not exist) plus a
\emph{proposal token}: a hash of the operation, the canonicalised terms, and
the current root. \cmd{put} recomputes the token; if the store has moved or a
spelling resolved differently, it refuses and asks the agent to propose again.

This matters more for agents than for humans. An agent may spend many seconds
of reasoning between reading and writing, during which the store can move; and
it cannot notice a surprise the way a person glancing at a terminal would.
Making the surprise into a refusal is the difference between a race and an
error message.

\paragraph{Every write is a speech act.} A successful write commits two
things: the knowledge change, and one signed provenance record per claim,
recording the author's key, the claim (a canonical
$\mathit{sub} \dsub \mathit{sup}$ pair), and the root the author was looking
at when they decided. Removals emit retraction records: retracting is also
something somebody did.

Notice where the signature goes. It is in a \emph{separate} store with its own
root, never inside the knowledge records --- because if authorship were part
of the hashed knowledge, two people who know the same thing would get
different roots, and G1 would be gone. \emph{Who said it} and \emph{what is
said} are deliberately different objects with different addresses.

\subsection{The review problem}
\label{sec:agents-review}

Here is the problem that will dominate agent-written knowledge, and it is not
a technical one:

\begin{keyidea}
Agents write faster than anyone can review. A store that accepts machine
writes at machine speed accumulates claims nobody has checked, and the value
of the store degrades to the value of its least careful writer.
\end{keyidea}

\odag{} offers three partial answers and one promising direction.

\paragraph{Attribution makes the problem tractable rather than solving it.}
Because every claim is signed, a reader can filter: accept claims from these
keys, ignore the rest. Standing is computed from verified records only, and
retraction is sticky per author. Forged records are listed but never counted.
This does not check anything; it makes \emph{whose} unchecked claims you are
carrying an explicit decision.

\paragraph{Endorsement is a separate speech act.} A reviewer can endorse a
claim without touching the knowledge graph. Reviewing is thus recorded work
that accumulates, rather than an event that happens once in a chat window and
is lost.

\paragraph{Disagreement is queryable structure.} If disjointness claims are
ever adopted, an inconsistency arriving by merge is not an error condition; it
is a non-empty \cmd{get(Cat, Dog)}. The consistency check is the query
primitive again --- which means finding the mess an agent made is the same
operation as using the store.

\paragraph{The direction: bound and order the review.}
Section~\ref{sec:learn-exploration} argued that attribute exploration is the
right shape here, and the reason belongs in this section. Exploration inverts
the flow: instead of the agent proposing arbitrary edits which a human must
review in arbitrary order, the \emph{system} computes the next question that
is not already answered, and the agent answers it. The review queue becomes
minimal (no redundant questions), ordered (each maximally informative given
the answers so far), and terminating. That is the difference between review
that scales and review that does not.

\subsection{Agents sharing memory}

The multi-agent case is where the properties compound, and it is worth a
concrete scenario.

Two agents work on the same domain on different machines, offline from each
other. Each files what it learns. Later they exchange roots.

\begin{enumerate}[leftmargin=1.7em]
\item If the roots are equal, they know the same things. \emph{One comparison,
  32 bytes, no data transferred} --- and because canonicality is semantic, this
  works even if they built their stores in different orders, by different
  routes, spelling values differently.
\item If the roots differ, a structural diff localises the disagreement to the
  exact differing records, without downloading either store wholesale.
\item Merging converges: each folds in the other's root and both reach the
  byte-identical result, whatever order the exchange happened in (G5). No
  server, no consensus, no lock.
\item Afterwards, each can ask ``which of these claims came from the other
  agent, and do I trust that key?'' --- because the provenance store merged
  too, conflict-free, by the same union rule.
\end{enumerate}

Compare the same scenario with two agents holding vector stores, or two
fine-tuned models: there is no equality test, no diff, no merge, and no
attribution. There is only ``ask them both and hope''.

\subsection{Token economics, briefly}

A practical point that matters at scale. An exact query returns exactly the
relevant items and a count. An approximate retrieval returns $k$ chunks chosen
by similarity, of which some are relevant, and the agent pays context tokens
for all of them and then reasons about which to believe. For queries that are
genuinely conjunctive --- ``refundable, in Japan, booked before March'' ---
the exact path is both cheaper and correct, and the cost difference grows with
the store.

The same applies to structure. An agent that needs to know what a store
contains can read a \cmd{about} record of a few hundred tokens rather than
sampling documents. Discoverability is a token-efficiency feature as much as a
correctness one.

\subsection{What \odag{} does not solve}

An honest list, because the value of the guarantees depends on not
overstating them.

\begin{itemize}[leftmargin=1.4em]
\item \textbf{Proofs attest structure, never truth.} Every certificate proves
  what was asserted and what follows from it --- never that the assertion is
  true of the world. That is the oracle problem, and no data structure solves
  it. The available answers are social (attribution and endorsement) and
  economic (Section~\ref{sec:crypto-bonds}).
\item \textbf{Classification is not automatic.} \odag{} has no defined
  concepts (\S\ref{sec:meet-trap}), so nothing places an item in the right
  category for you. That job falls to the agent --- which is a reasonable
  division of labour in 2026, but it means the quality of the store depends on
  the quality of the classifier.
\item \textbf{Names must be agreed.} Two agents using different words for the
  same idea produce two nodes, and no mechanism notices. This is the price of
  Section~\ref{sec:keep-names}, and it is real.
\item \textbf{No relations.} ``Alice authored Doc1'' cannot be said. Agents
  will try; the surface logs what they fail to express, which is how the
  system will learn whether the wall needs moving.
\item \textbf{Retraction does not travel.} A peer who still holds a removed
  edge re-adds it on merge. Single-writer and publish-to-readers deployments
  keep removals; concurrent writers do not.
\end{itemize}
